\documentclass[12pt,a4paper]{article}

% Page formatting
\usepackage[margin=25mm]{geometry}
\usepackage{setspace}
\onehalfspacing

% Font and encoding
\usepackage[T1]{fontenc}
\usepackage[utf8]{inputenc}
\usepackage{lmodern}  % Clean, scalable font

% Section formatting
\usepackage{titlesec}
\titleformat{\section}{\large\bfseries}{\thesection}{1em}{}

% Header/footer
\usepackage{fancyhdr}
\pagestyle{fancy}
\fancyhf{}
\renewcommand{\headrulewidth}{0pt}
\fancyfoot[C]{\thepage}

% Hyperlinks
\usepackage[hidelinks]{hyperref}

% Figures and graphics
\usepackage{graphicx}
\usepackage{caption}
\usepackage{subcaption}

\usepackage{amsmath,amssymb}

\usepackage[backend=biber,style=numeric]{biblatex}
\addbibresource{references.bib}

\title{\vspace{-2cm}Technical Abstract: \textit{Acoustics of Quadcopter Propellers}}
\author{Louis Pender \\
        Lucy Cavendish College, University of Cambridge}
\date{\vspace{-5mm}}

\usepackage{titling}
\setlength{\droptitle}{-1.5cm}

\begin{document}

\maketitle

\thispagestyle{empty} % No header/footer on title page

\section*{Abstract}

\newpage
\tableofcontents

\newpage
\section{Introduction}

\subsection{Background}

\subsection{Motivation}

\subsection{Objectives}

\section{Theory}
% sentence about the theory of this report not going into a lot of detail because
% of the broad scope of the report but will be crucial nevertheless
The relevant theory in this report covers acoustic and aerodynamic

% aero side

% explain the origin of thrust and torque
% dependence of these aeroydynamic forces with speed, diameter, introduce invarient coefficients
% its more complicated than that, BET: nonlinear version of airfoil theory
% justify why its more complicated than that. we dont want to overcomplicate things

% derivation of nonlinear solver used to predict aerodynamic performance of propellers
% extension to swept blades

% explain that this doesnt include swirl or tip losses where the blade element theory breaks down

% acoustic side

% explain origin of aerodynamic noise, maybe an acoustic analogy
% introduce intuitively how the sound pressure level would be expected to vary with speed, diameter, distance
% justify the focus on lift being the primary source of propeller noise
% introduce concepts of interference with definition of wavenumber to explain directivity.

% explain model of dipole point sources (simplest model )
% explain why the integral of many helically convected sources models the propeller
% explain how changing location of sources can cause destructive interference for some wavenumbers

% implementation of the model, extension to looped propellers

\section{Methodology}

% brief explaination of the overall approach

\subsection{Experimental}
% justification for separating acoustic and aerodynamic measurements
% where was the testing done

\subsubsection{Aerodynamic setup}
% show the setup used and explain the components
% 200g, 1Kg load cells, hx711@10Hz pi pico. Odrive S1, AS5047p encoder, 10kOhm motor thermistor
% explain some of the choices made in the setup
% explaination of calibration procedure
% control of variables 

\subsubsection{Acoustic setup}
% show the setup used and explain the components
% explain some of the choices made in the setup
% control of variables


\subsubsection{Data Collection and Processing}
% software used, majority python, C++
% key libraries used numpy, scipy, matplotlib and xfoil

% rms normalised hanning window to prevent spectral leakage
% frequency domain interpolation of microphone calibration data

% method of integrating harmonics
% log regression of distance and speed

\subsection{Theoretical}



\section{Results}

\section{Discussion}

% specific results of the two setups

\section{Conclusion}

\subsection{Improvements}


% better experimental setup
%% 
%% 

% better theoretical models
%% accounting for swirl and tip losses in current solver
%% interpolating airfoil data with reynolds number in current solver
%% panel method, wake vortex model
%% inclusion of volume displacement noise

\subsection{Future Work}
Alongside the improvements mentioned in the previous section, there are a number of avenues for future work.
% what questions were left unanswered


\end{document}
